\documentclass[10pt,a4paper,dvipdfmx]{jsarticle}
\usepackage[utf8]{inputenc}
\usepackage[top=15mm, bottom=15mm, left=15mm, right=15mm]{geometry}
\usepackage{graphicx}
\usepackage{xcolor}
\usepackage{fancyhdr}
\usepackage{titlesec}
\usepackage{calc}
\usepackage{ulem}
\usepackage{tcolorbox}
\usepackage{pgffor} % ノート罫線ループ用
\usepackage{otf}
\usepackage{multicol}
\newcounter{blanknum}
\tcbuselibrary{skins,breakable}

% ==============================================
% ★ 穴埋めモードの切り替え設定 ★
% 解答プリント（赤字）にする場合は \answertrue
% 生徒用プリント（空欄）にする場合は \answerfalse に変更
\newif\ifanswer
\answertrue % ←ここを変更
% ==============================================

% 穴埋め用マクロ（白文字化＋下線で文字幅を維持）
\newcommand{\aw}[1]{%
  \ifanswer
    %\textcolor{red}{\textbf{#1}}%
    \stepcounter{blanknum}%
    % 番号を左にはみ出させつつ、
    % そのぶんだけ左余白を確保
    \ifnum\value{blanknum}<10
      \hspace{1.2em}%
    \else
      \ifnum\value{blanknum}<100
        \hspace{1.6em}%
      \else
        \hspace{1.9em}%
      \fi
    \fi
    \llap{\textsubscript{\scriptsize（\theblanknum）}}%
    \uline{\textcolor{red}{\textbf{\Large{#1}}}}%
  \else
    \stepcounter{blanknum}%
    % 番号を左にはみ出させつつ、
    % そのぶんだけ左余白を確保
    \ifnum\value{blanknum}<10
      \hspace{1.2em}%
    \else
      \ifnum\value{blanknum}<100
        \hspace{1.6em}%
      \else
        \hspace{1.9em}%
      \fi
    \fi
    \llap{\textsubscript{\scriptsize（\theblanknum）}}%
    \uline{\textcolor{white}{\textbf{\Large{#1}}}}%
  \fi
}

% 画像挿入用マクロ
\newcommand{\insertimage}[2]{%
  \begin{center}
    \tcbox[colback=gray!5,colframe=gray!50,boxrule=1pt,arc=2pt,boxsep=5pt,width=0.9\linewidth]{%
      \begin{tabular}{c}
        \textbf{【画像挿入指定箇所】 #1} \\
        \footnotesize ※ここに \texttt{\textbackslash includegraphics[width=15cm]\{#2\}} を挿入してください。
      \end{tabular}
    }
  \end{center}
}

% 補遺（深堀りコラム）用ボックス
\newtcolorbox{hocho}[1][]{
  enhanced, breakable,
  colback=gray!5, colframe=darkgray,
  fonttitle=\bfseries,
  title={【補遺】 #1},
  attach boxed title to top left={yshift=-2mm, xshift=5mm},
  boxed title style={colback=darkgray, colframe=darkgray}
}

% ==============================================
% ★ ノート（罫線）ページ自動生成マクロ ★
% ==============================================
% 手動で1ページ分のノート（MEMO）を挿入するコマンド
\newcommand{\insertnotepage}{
  \clearpage
  \thispagestyle{fancy}
  \fancyhead[L]{\textbf{MEMO}}
  \fancyhead[R]{}
  \vspace*{2mm}
  \foreach \i in {1,...,20} {
    \noindent\textcolor{gray!40}{\rule{\textwidth}{0.4pt}}\par\vspace{8mm}
  }
  \clearpage
}

% ページ合わせ自動化コマンド：両面印刷時、裏面が白紙になる場合のみノートを自動挿入
\newcommand{\fillnotepage}{
  \clearpage
  \ifodd\value{page}
    % 現在奇数ページなら、前のページ（偶数）で終わっているので何もしない
  \else
    % 現在偶数ページなら、裏面が白紙になってしまうのでノートページにする
    \insertnotepage
  \fi
}
% ==============================================

% ヘッダー・フッター設定
\pagestyle{fancy}
\fancyhf{}
\fancyhead[L]{\textbf{高３世界史 一学期中間試験対策 補助教材}}
\fancyhead[R]{\textbf{戦後のアジア・第三世界}}
\fancyfoot[C]{\thepage}
\renewcommand{\headrulewidth}{0.4pt}

% セクションタイトルのデザイン
\titleformat{\section}[hang]
  {\Large\bfseries}{\colorbox{darkgray}{\color{white}\thesection}}{1em}{\uline}
\titleformat{\subsection}[hang]
  {\large\bfseries}{\colorbox{gray!30}{\thesubsection}}{1em}{}
\titleformat{\subsubsection}[hang]
  {\normalsize\bfseries}{}{0em}{■ }

\begin{document}

\begin{center}
  {\Huge \textbf{戦後のアジア・第三世界}}\\[0.5em]
  {\Large 高３世界史　１学期中間}
\end{center}
\vspace{1em}

% ======================================================================
\section{第二次世界大戦後の植民地体制の崩壊と超大国の思惑}
% ======================================================================

\subsection{植民地独立の本格化}
第二次世界大戦後、アジア、アフリカの植民地各地で粘り強い独立運動が爆発した。これには二つの大きな要因がある。
第一に、第2次世界大戦による総力戦の結果、\aw{イギリス}や\aw{フランス}といった旧宗主国が経済的にも軍事的にも著しく疲弊し、地球規模の植民地を維持することが物理的に困難になったことである。
第二に、戦後の国際政治を主導した\aw{アメリカ}と\aw{ソ連}という二つの超大国が、旧来の西欧型帝国主義による植民地体制の維持には無関心であり、むしろ植民地の「\aw{解放}」に共感を示したことである。ただし、この「共感」は冷戦下の激しい覇権争いの中での極めて戦略的な打算に基づいていた。

\begin{itemize}
    \item \textbf{アメリカの思惑}：独立運動が\aw{反共}（反共産主義）につながり、西側陣営の市場と軍事同盟に組み込める場合は積極的に独立を支援した。一方で、共産主義者が独立運動の主導権を握る場合は、独立を徹底的に妨害し軍事介入も辞さなかった（例：ベトナム戦争）。
    \item \textbf{ソ連の思惑}：帝国主義陣営の市場と資源供給地を奪い、社会主義圏を拡大させるという大戦略のもと、各国の民族解放戦線などの武力闘争を物心両面で支援した。
\end{itemize}

\begin{hocho}[第二次世界大戦の「複合的な性格」と三つの道]
第二次世界大戦は単なる帝国主義間の領土争いではなく、「ファシズム対反ファシズム」というイデオロギー戦争の性格も持っていた。そのため、植民地の指導者たちは独立を勝ち取るために複雑な選択を迫られた。
(1)\textbf{反ファシズムの立場から宗主国へ戦争協力}：各国の\aw{共産党}は概ねこの道を選んだ。ファシズム（日・独・伊）を倒すため、一時的に宗主国（英・仏・蘭）と協力し、戦後の譲歩を引き出そうとした。
(2)\textbf{独立を重視し、宗主国と戦う日本と手を結ぶ}：東南アジアの「\aw{ナショナリスト}」に多く見られた（例：インドネシアのスカルノ、ビルマのアウン・サン）。「敵の敵は味方」の論理で日本軍から軍事訓練を受け、独立軍の基盤を築いた。
(3)\textbf{宗主国にもファシズム勢力にも協力しない}：インドの\aw{国民会議派}がこの困難な選択をした。「インドから立ち去れ」運動を展開し、ファシズムには反対しつつもイギリスの戦争遂行には非協力を貫いた。
\end{hocho}

\subsection{宗主国の妥協的再編と帝国主義の末路}
イギリスとフランスは、独立が不可避であることを悟りつつも、旧植民地への経済的・政治的特権を残すために緩やかな国家連合への再編を試みた。
\begin{itemize}
    \item \textbf{イギリス}：大英帝国からイギリス連邦（ブリティッシュ・コモンウェルス）へと移行していた体制を、さらに対等な主権国家の連合であることを強調する\aw{コモンウェルス・オブ・ネイションズ}へと再編した。※イギリスへの強い恨みを持つアイルランドは1948年に共和国法を可決し、1949年に完全に脱退した。
    \item \textbf{フランス}：第4共和制下で\aw{フランス連合}を形成したが、インドシナ戦争の敗北で破綻。1958年、アルジェリア危機のさなかに成立した第5共和制の下で、\aw{フランス共同体}（コミュノテ・フランセーズ）へと移行。本国との経済関係を維持したままの独立を可能にしたためアフリカでの独立が急増したが、離脱する国が相次ぎ消滅。1970年代以降は、フランス語を共通言語とする文化的紐帯である\aw{フランコフォニー}へと姿を変えた。
\end{itemize}

% ======================================================================
\section{アジアにおける「民族解放」の連鎖}
% ======================================================================

1940年代後半から50年代にかけて、東南アジアから中東に至るまで、民族解放の嵐が吹き荒れた。

\subsection{東南アジアの独立の軌跡}
\begin{itemize}
    \item \textbf{インドネシア}：大日本帝国の降伏直後である1945年8月17日、\aw{スカルノ}を大統領として独立を宣言。しかし、再支配を狙うオランダが軍事侵攻し、4年にわたる過酷な\aw{インドネシア独立戦争}が勃発した。ゲリラの頑強な抵抗と、オランダの残虐行為に対するアメリカからの経済援助停止の圧力により、1949年にオランダが主権委譲を承認。当初は連邦共和国（オランダ側の分割統治の罠）であったが、翌1950年に単一国家の\aw{インドネシア共和国}へと再編した。
    \item \textbf{ベトナム}：1945年の八月革命により、ホー・チ・ミン主導で独立を宣言するも、フランスがこれを認めず\aw{インドシナ戦争}（1946〜54）へ突入。冷戦下での「熱戦」を経験することとなる。
    \item \textbf{ラオス、カンボジア}：当初はフランス連合内での不完全な王国独立であったが、1953〜54年に完全な主権承認を得た。しかし内戦を経て、ラオスは1975年に社会主義の\aw{ラオス人民民主共和国}に、カンボジアは1993年に\aw{カンボジア王国}となった。
    \item \textbf{フィリピン}：アメリカは戦前の1934年の法律で10年後の独立を約束しており、戦後の1946年に\aw{フィリピン共和国}として独立。ただし、米軍基地の存続などアメリカへの強い従属が続いた。
    \item \textbf{マレーシアとシンガポール}：1957年、イギリス連邦内の自治領として\aw{マラヤ連邦}が独立。1963年にボルネオ島北部などと合体し\aw{マレーシア連邦}が成立した。しかし、マレー人優遇政策と対立した華人主体の\aw{シンガポール}が、1965年に分離独立（事実上の追放）した。初代首相は\aw{リー＝クアンユー（李光耀）}である。
\end{itemize}

\begin{hocho}[シンガポール「涙の独立宣言」]
1965年8月9日、シンガポールの初代首相リー・クアンユーはテレビカメラの前で涙を流しながら独立を宣言した。資源も水もなく、国土は東京23区ほどの狭小な島であるシンガポールにとって、マレーシアからの分離は「死の宣告」に等しかったからだ。「私は生涯のすべてをマレーシアとの統合に捧げてきた…」と語り落涙した彼は、この絶望的な状況から、強権的な開発独裁と徹底した外資誘致、教育改革によって、わずか一代でシンガポールをアジア屈指の豊かな先進国へと押し上げた。
\end{hocho}

\subsection{インド・パキスタンの分離独立と中東の再編}
\begin{itemize}
    \item \textbf{インドの分断}：イギリスの「分割統治」の遺恨により、1947年、ヒンドゥー教徒中心の\aw{インド}と、イスラーム教徒中心の\aw{パキスタン}が分離独立。両国は帰属を巡って\aw{カシミール問題}を引き起こし、三度にわたるインド・パキスタン戦争を戦った。
    \item \textbf{バングラデシュ独立}：パキスタンはインドを挟んで東西に分断されており、ウルドゥー語を強制し政治・経済を独占する西パキスタンに対し、1971年に東パキスタンが反乱。インドの軍事支援を受け、\aw{バングラデシュ}として独立した。
    \item \textbf{中東の国家形成}：第一次大戦後のサン・レモ会議で英仏の委任統治領とされた地域から、1943年に\aw{レバノン}、1946年に\aw{シリア}（一時はエジプトと合邦しアラブ連合共和国に）とトランスヨルダン（1950年に\aw{ヨルダン}＝ハシェミット王国と改称）が独立。1948年にはユダヤ人による\aw{イスラエル}が独立宣言し、中東戦争の火蓋が切られた。
\end{itemize}

\begin{center}
 \includegraphics[width=15cm]{asiamapfull.png}
\end{center}

% ======================================================================
\section{アフリカの独立運動とアラブ・民族主義の高揚}
% ======================================================================

1950年代から60年代にかけて、アフリカ大陸の地図は劇的に塗り替えられた。

\subsection{北アフリカ・アラブ諸国の独立とエジプト革命}
1950年代前半から、リビア（1951年）、チュニジア・モロッコ・スーダン（1956年）など北アフリカのアラブ諸国が相次いで独立を果たした。その精神的支柱となったのがエジプトの動向である。
\begin{itemize}
    \item \textbf{エジプト革命（1952年）}：イギリスに追従し腐敗していた王政（\aw{アリー朝}）に対し、軍内部の秘密結社「\aw{自由将校団}」がクーデタを起こして国王を追放し、共和政を樹立した。初代大統領ナギブを経て、実権を握った第2代大統領\aw{ナセル}がアラブ民族主義の英雄となる。
    \item \textbf{スエズ戦争（1956年）}：ナセルはナイル川の治水と近代化のため\aw{アスワン＝ハイダム}の建設を計画したが、アメリカなどが資金援助を撤回。これに対抗し、1956年7月に\aw{スエズ運河国有化宣言}を断行した。
    これに激怒したイギリス・フランスは、\aw{イスラエル}と同調してエジプトを軍事攻撃（スエズ戦争／第2次中東戦争）。しかし、旧態依然とした帝国主義の横暴に対し、ソ連とアメリカを含む国際世論はエジプトに同情し、英仏に撤退を強要した。これにより中東におけるイギリスの影響力は完全に失墜し、エジプトは自らの大国化をはかり、\aw{アラブ民族主義}の熱狂的な高揚をもたらした。
    \item \textbf{アルジェリア独立戦争（1954-62年）}：100万人以上のフランス人入植者（コロン）がいたアルジェリアでは、\aw{アルジェリア民族解放戦線（FLN）}主導の過酷な独立戦争が勃発。フランス本国でも独立支持派と軍部・入植者側で内戦寸前の対立が起き、第四共和政が崩壊。新たに第五共和政を樹立したド＝ゴールが、1962年の\aw{エビアン協定}でついに独立を承認した。
\end{itemize}

\begin{hocho}[フランス軍の拷問と「アルジェリアの戦い」]
アルジェリア独立戦争は、双方がテロと残虐行為を繰り返す凄惨な泥沼の戦いであった。首都アルジェのカスバ（旧市街）に潜むFLNのゲリラに対し、フランス軍の落下傘部隊は徹底した掃討作戦を行い、水責めや電気ショックなどの過酷な拷問を用いて自白を強要した。この歴史の闇は、映画『アルジェの戦い』（1966年）でドキュメンタリー・タッチで生々しく描かれ、後の世界中のゲリラ戦の教科書になると同時に、フランスにとって長く語ることをタブーとされる「汚い戦争」の記憶となった。
\end{hocho}

\subsection{サブサハラ・アフリカの独立の波}
1950年代後半から、サハラ砂漠以南（サブサハラ）のアフリカでも独立の波が押し寄せた。
\begin{itemize}
    \item \textbf{ガーナの独立}：1957年、サハラ以南のブラックアフリカで最初に独立を果たしたのが\aw{ガーナ}（旧英領ゴールドコースト）である。初代首相（のち大統領）の\aw{ンクルマ}は、「アフリカの独立は、全アフリカの解放と結びつかない限り無意味である」と語り、汎アフリカ主義に基づく「\aw{アフリカ合衆国}」構想を掲げた。
    \item \textbf{アフリカの年（1960年）}：フランス領のアフリカ諸国を中心に、一挙に17カ国が独立を果たし「アフリカの年」と呼ばれた。国連総会でも植民地主義への反対決議（植民地独立付与宣言）が採択され、脱植民地化が加速した。ただし、新興独立国の多くは、植民地政府の「\aw{収奪機能}」を一部の現地エリートが継承するという不安定な構造的欠陥を抱えていた。
    \item \textbf{アフリカ統一機構（OAU）}：1963年、エチオピアのアディスアベバでアフリカ独立諸国首脳会議が開かれ、\aw{パン＝アフリカニズム}を理念とするアフリカ統一機構（OAU）が創設された。これは2002年に、より強力な経済的・政治的統合を目指す\aw{アフリカ連合（AU）}へと発展的に解消された。
\end{itemize}

\subsection{ポルトガル帝国の崩壊と脱植民地化の完了}
最も遅くまで植民地支配に固執した最後の植民地帝国ポルトガルは、泥沼の植民地戦争で疲弊し、1974年に本国で起きた\aw{ポルトガル革命}（カーネーション革命・軍事独裁崩壊）により方針を転換。1975年にアンゴラ、モザンビーク、ギニアビサウなどが次々と独立した。
さらに、1990年に南アフリカの不法占領下にあった\aw{ナミビア}が独立したことで、アフリカ大陸の政治的な脱植民地化はほぼ完了した。
南アフリカ共和国本国でも、国際社会の強力な経済制裁と国内の抵抗運動により、1991年に\aw{黒人隔離諸法}（アパルトヘイト）の撤廃が宣言され、1994年の全人種参加の民主選挙によりネルソン・マンデラ大統領が誕生し、制度的な人種差別は完全消滅した。

\begin{center}
  \includegraphics[width=15cm]{africamapfull.png}
\end{center}

% ======================================================================
\section{大日本帝国の崩壊とアジアの「熱戦」の始まり}
% ======================================================================

\subsection{権力の真空と冷戦の最前線化}
1945年8月、大日本帝国の敗戦は東アジア・東南アジアに巨大な「権力の真空」を生み出した。ヨーロッパでは冷戦は「冷たい」対立（鉄のカーテン）にとどまっていたが、国境線が未確定であったアジアにおいては、大国間の代理戦争として直接火を噴く「\aw{熱戦}」となった。

\subsection{中国革命の決着と分断}
日本では終戦を迎えた直後から、中国大陸では\aw{\UTF{8523}介石}率いる国民党（アメリカ支援）と、\aw{毛沢東}率いる共産党（ソ連支援）による第2次\aw{国共内戦}が再開した。
農地改革の約束で広範な農民の支持を集め、日本軍が残した武器を接収した共産党軍（人民解放軍）が最終的に勝利し、1949年10月に\aw{中華人民共和国}が成立した。敗れた国民党勢力は、艦艇や黄金の準備金とともに\aw{台湾}へ逃れ、中華民国の体制を維持した。

\subsection{朝鮮半島の分断と朝鮮戦争}
日本の過酷な植民地支配から解放された朝鮮半島では、直ちに建国準備委員会が「朝鮮人民共和国」の成立を宣言した。しかし、反共の防波堤を構築したいアメリカはこれを一切認めず、ソ連との間で北緯38度線を境界線とする「分割占領」を行い、国際信託統治を想定した米ソ交渉が行われたが難航した。
最終的にアメリカは国連を動かし、南半部のみの単独選挙を強行。1948年8月〜9月、南に李承晩を大統領とする\aw{大韓民国}、北に金日成を首相とする\aw{朝鮮民主主義人民共和国}が成立し、民族の分断が決定づけられた。

\begin{itemize}
    \item \textbf{開戦と初期の攻防}：1950年6月25日早朝、ソ連のスターリンの暗黙の了解を得た北朝鮮軍が北緯38度線を一斉に越えて南進し、\aw{朝鮮戦争}（1950-53年）が勃発した。油断していた韓国軍は総崩れとなり、同年8月には半島南端の釜山橋頭堡まで追い詰められた（最南進戦線）。
    \item \textbf{国連軍の反撃と中国の介入}：アメリカはマッカーサー元帥を総司令官とする国連軍（実質は米軍主体の多国籍軍）を組織。1950年9月14日、奇跡的な「仁川上陸作戦」を成功させて背後を突き、ソウルを奪回。一気に北進し、11月には中朝国境の鴨緑江にまで迫った（最北進戦線）。しかし、「唇亡びれば歯寒し」と自国の安全保障に脅威を感じた中国が、彭徳懐率いる約30万人の「\aw{人民義勇軍}」を大挙派遣。圧倒的な人海戦術と山岳ゲリラ戦で国連軍を押し返し、1951年3月には再びソウル以南まで進出した（最南下線）。
    \item \textbf{膠着と休戦}：その後、北緯38度線付近で陣地戦（高地戦）の泥沼に陥った。1951年6月から板門店などで休戦交渉が始まったが、捕虜の強制送還を巡って長引いた。1953年3月にソ連の独裁者\aw{スターリンが死去}したことでようやく東側陣営の態度が軟化し、同年7月に休戦協定が締結された。※あくまで「休戦」であり、\aw{講和条約}は現在に至るまで結ばれていない。
\end{itemize}

\begin{hocho}[なぜ国連軍にソ連は「拒否権」を使えなかったのか？]
朝鮮戦争への介入を決める国連安全保障理事会において、なぜ常任理事国であるソ連は拒否権を発動しなかったのか。実は当時、ソ連は安保理を「ボイコット（欠席）」していたのである。国連の中国代表権が中華人民共和国（北京）ではなく中華民国（台湾）に与えられ続けていることに抗議するためだった。この欠席の隙を突いてアメリカは軍事介入の決議を強行し、「国連軍」の御旗を手に入れた。この大失態以降、ソ連は二度と安保理の重要な会議をボイコットすることはなくなった。
\end{hocho}

\begin{center}
   \includegraphics[width=15cm]{koreanwarmap.png}
\end{center}

% ======================================================================
\section{朝鮮半島の戦後復興と開発独裁の台頭}
% ======================================================================

戦火で焦土と化した朝鮮半島の南北両国は、東西冷戦のショーウィンドウとして、それぞれの陣営から大規模な援助を受けて復興の道を歩んだ。

\subsection{北朝鮮の社会主義的工業化}
ソ連、中国、東欧など社会主義陣営からの莫大な経済・技術援助を受けた朝鮮民主主義人民共和国は、1960年にはいち早く「工業国」化を成し遂げた。
歴史的に見て、日本統治時代に豊富な水資源と鉱産資源を活用して北部に巨大な水力発電所や化学コンビナートが建設されていたため（これを「\aw{南農北工}」と呼ぶ）、北朝鮮の重工業化は韓国よりも先行していた。政治面では、朝鮮労働党の指導のもと農業の集団化と工業の国有化が進み、\aw{金日成}は政敵（延安派やソ連派）を徹底的に粛清し、絶対的な権力集中（主体思想）を確立した。

\subsection{大韓民国の混乱と「四月革命」}
一方の大韓民国は、アメリカからの無償援助によって復興を図った。とくにアメリカの余剰農作物（小麦、砂糖、綿花）を受け入れて加工する製粉・精糖・紡織工業、いわゆる「\aw{三白産業}」が発達した。また、日本が残した莫大な財産（帰属財産）の払い下げを一部の特権的な資本家が独占し、彼らが後の独占財閥（チェボル）へと成長した。
しかし、アメリカからの大量の無償農産物の流入は国内の農産物価格を暴落させ、農地改革で生まれたばかりの自作農を困窮させた。土地を手放した農民はソウルなどの都市部へスラムを形成しつつ流入し、安価な工場労働の担い手となった。
1950年代後半、アメリカの援助が削減され有償化に切り替わると経済は深刻な不況に陥った。初代大統領\aw{李承晩}（イ・スンマン）は強権を用いて1960年3月の大統領選で四選を果たすが、大規模な不正選挙であったため学生や市民の怒りが爆発。激しい抗議デモ（警官隊の発砲で多数の死傷者）の末、4月に李承晩は退陣しハワイへ亡命した。これを\aw{四月革命}と呼ぶ。その後、議院内閣制の新憲法下で張勉政権が成立したが、与党内紛で弱体化し社会の混乱は収まらなかった。

\subsection{朴正煕の登場と「漢江の奇跡」}
1961年5月、社会の混乱を力で抑え込むべく、\aw{朴正煕}（パク・チョンヒ）ら一部の将校が軍事クーデタを起こし、議会を解散させて実権を掌握した。アメリカのケネディ政権は当初難色を示したが、朴が強烈な「反共」を掲げたため事実上これを黙認した。

朴政権の至上命題は「経済の自立化」であった。投資のための外貨を喉から手が出るほど欲していた朴は、国内の猛烈な反対デモを戒厳令で弾圧しながら、1965年に日本との国交を回復する\aw{日韓基本条約}を締結した。これにより日本から無償3億ドル・有償2億ドルの莫大な経済協力金を引き出した。さらに1964年からは、アメリカの要請に応えて\aw{ベトナム戦争への派兵}（延べ約32万人）を強行し、軍需物資の調達や送金による莫大な特需（ベトナム特需）を得た。
これらの血みどろの財源を基に、高速道路や浦項製鉄所などの重化学工業化を国主導で推進し、ソウルを流れる川の名をとって「\aw{漢江の奇跡}」と呼ばれる驚異的な高度経済成長を実現した。しかしこれは、労働運動を過酷に弾圧し、反対派をKCIA（中央情報部）を用いて粛清する強権的な「\aw{開発独裁}」の産物であった。

\begin{hocho}[「請求権の完全かつ最終的な解決」と韓国のジレンマ]
1965年の日韓基本条約締結時、日本政府は過去の植民地支配への「賠償」ではなく、あくまで独立祝いの「経済協力金」という名目で資金を供与した。この際、両国間で「請求権は完全かつ最終的に解決された」とする協定が結ばれた。朴正煕政権は、この莫大な資金を個人の被害者への補償にはほとんど回さず、道路やダム、製鉄所などの国家インフラ建設に全額投入してしまったのである（これが漢江の奇跡の原資である）。しかし、この「国家が個人の請求権を犠牲にして経済発展を優先した」という歴史的経緯が、後の慰安婦問題や徴用工問題など、現在に至る日韓の歴史認識の火種を残すこととなった。
\end{hocho}

% ======================================================================
\section{ベトナム戦争：独立の総力戦から泥沼の代理戦争へ}
% ======================================================================

\subsection{インドシナ戦争とディエンビエンフーの屈辱}
1945年8〜9月の\aw{八月革命}により、独立の父\aw{ホー＝チ＝ミン}主導でベトナム民主共和国が独立宣言を行った（社会主義政権）。しかし、第二次大戦の勝戦国として植民地帝国の復活を目論むフランスはこれを武力で潰そうとし、1946年から\aw{インドシナ戦争}（1946-54年）が勃発した。
フランスは1949年、ベトナム人のナショナリズムを取り込むため、阮朝最後の皇帝\aw{バオ＝ダイ}を国家主席とする傀儡政権「ベトナム国」を樹立した。しかし、1954年の山岳要塞における\aw{ディエンビエンフーの戦い}で、中国の支援を受けたベトナム人民軍が、空挺部隊の精鋭を揃えたフランス軍を包囲・殲滅し、大敗北を与えた。
白人帝国主義軍がアジアの民族解放軍に真正面から敗北したこの歴史的事件により、同年7月に\aw{ジュネーヴ協定}が結ばれた。\aw{北緯17度線}を軍事境界線として一時的に南北を分割し、2年後に南北統一選挙を実施して完全独立を図ることが約束された。

\subsection{アメリカの介入とゴ・ディン・ジエムの独裁}
しかし、冷戦下のアメリカ（アイゼンハウアー政権）は、「ベトナムが共産化すれば、ドミノ倒しのように東南アジア全域が赤化する」というドミノ理論の強迫観念に囚われていた。統一選挙を行えば国民的英雄ホー・チ・ミンが圧勝することは明白であったため、アメリカは選挙を拒否。1955年にバオ＝ダイを追放し、反共主義者でカトリック教徒の\aw{ゴ＝ディン＝ジエム}を大統領とする\aw{ベトナム共和国}（南ベトナム）を強引に樹立させた。
ジエム政権は、アメリカからの莫大な援助を背景に、仏教徒（国民の大多数）の弾圧や縁故主義による腐敗した独裁を敷いた。この暴政への反発から、1960年に南部の農民や知識人が結集して\aw{南ベトナム解放民族戦線（NLF）}（通称ベトコン）が組織され、北ベトナム（ベトナム労働党）の支援を受けてゲリラ戦を開始した。
ケネディ政権は軍事顧問団を増派しつつ内政改革を求めたが、ジエムが拒絶したため、1963年11月に南ベトナム軍部のクーデタを黙認し、ジエムは暗殺された。

\subsection{「北爆」開始と総力戦への転化}
ジエム暗殺後も南ベトナムの混乱は続き、解放戦線が勢力を拡大した。これに焦ったジョンソン政権は、1964年8月の\aw{トンキン湾事件}（米駆逐艦が北ベトナムの魚雷艇に攻撃されたとする事件。後に米軍の挑発・捏造であったことが暴露される）を口実に、議会の白紙委任を得て北ベトナムへの報復爆撃を開始した。
1965年2月からは絨毯爆撃である「\aw{北爆}」を本格化し、7月には海兵隊などの地上部隊を大規模に派遣。アメリカが直接戦闘の泥沼に足を踏み入れる本格的な「ベトナム戦争」が始まった。

この戦争の決定的な勝敗の要因は「モチベーションの非対称性」にあった。ベトナム人にとっては、相手がフランスからアメリカという巨大な帝国に変わっただけの、民族の存亡をかけた「\aw{独立のための総力戦}」であり、どんな犠牲（200万人以上の死者）を払っても戦い抜く覚悟（フルパワー）があった。一方、アメリカにとっては、ソ連や中国との全面核戦争を避けつつ共産主義の拡大を防ぐという冷戦の「\aw{限定戦争}」に過ぎず、明確な勝利条件もないまま若者がジャングルで命を落とす状況に、次第に士気と大義（やる気）を失っていったのである。

\begin{center}
   \includegraphics[width=8cm]{vietnamwarmap.png} 
   \includegraphics[width=8cm]{vnmw2.png} 
\end{center}

\subsection{テト攻勢と泥沼からの撤退}
北ベトナムは、1950年代後半から表面化していた\aw{中ソ対立}を巧みに利用し、「社会主義の盟主」を競い合うソ連と中国の両方から最新兵器（ミグ戦闘機や地対空ミサイルなど）の援助を大量に引き出した。一方、アメリカはイギリスやフランスなど西側の同盟国に派兵を求めたが拒否され（韓国などが派兵）、第三世界からは激しく批判され国際的に孤立していった。

1968年1月末、ベトナムの\aw{旧正月（テト）}の休戦期間を突き、解放民族戦線と北ベトナム正規軍が南部の主要都市数十カ所で一斉蜂起する「\aw{テト攻勢}」を敢行した。アメリカの聖域であったサイゴンのアメリカ大使館が一時ゲリラに占拠される映像は、テレビを通じて全米のお茶の間に届けられた。軍事的には米軍が圧倒的な火力でゲリラを掃討し勝利したものの、政治的・心理的には「もうすぐ戦争は終わる」と信じていたアメリカ国民に決定的な絶望を与え、反戦運動が爆発した。同年3月、ジョンソン大統領は北爆の部分的停止と、次期大統領選への不出馬を表明し事実上の敗北を認めた。

1969年に就任した\aw{ニクソン}大統領は「\aw{戦争のベトナム化}」（米軍を撤退させ、兵器を与えた南ベトナム軍に血を流させる）というニクソン・ドクトリンを発表。しかし一方で、北ベトナムの補給路（ホーチミン・ルート）を絶つため、隣国のカンボジアやラオスへ侵攻し、密かに絨毯爆撃を行ってインドシナ全域に戦火を拡大させた。

国内外の猛烈な批判と反戦運動に耐えきれず、1973年1月、ついに\aw{ベトナム（パリ）和平協定}が調印された。南北ベトナム、南ベトナム臨時革命政府（解放戦線）、アメリカの四者による協定で、米軍は完全撤退した。
アメリカという後ろ盾を失った南ベトナム軍は崩壊を続け、1975年3月に北ベトナム軍が大攻勢（春季大攻勢）を開始。4月30日、北ベトナムの戦車がサイゴンの大統領府の鉄柵を押し倒し、サイゴンが陥落した。
1976年、ついに南北は統一され、ハノイを首都とする\aw{ベトナム社会主義共和国}が成立した。

\begin{hocho}[テレビの戦争と「ソンミ村虐殺事件」]
ベトナム戦争は、史上初めて「テレビで生中継された戦争」であった。ジャーナリストが最前線に入り込み、検閲なしに血まみれの兵士や泣き叫ぶ子供の映像を世界に配信した。とくに世界を震撼させたのが1968年に発覚した「ソンミ村虐殺事件」である。ゲリラと一般住民の区別がつかない極度の恐怖と精神的疲労の中、米軍の小隊が無抵抗の老人、女性、子供を含む村民500人以上を機関銃で無差別虐殺した。この地獄の現実が雑誌にスクープされたことで、アメリカの掲げた「自由と民主主義の防衛」という大義名分は完全に崩壊し、反戦運動は学生だけでなく一般市民へも爆発的に広がったのである。
\end{hocho}

% ======================================================================
\section{中華人民共和国の成立と台湾の歩み}
% ======================================================================

\subsection{「反帝、反封建」の中国革命の結末}
1949年10月1日、天安門の楼上から\aw{毛沢東}が\aw{中華人民共和国}の成立を高らかに宣言した。これは、アヘン戦争以来100年以上にわたる「半植民地・半封建社会」からの解放（反帝、反封建の革命の最終的結末）を意味した。国家は\aw{中国共産党}によって指導され、毛沢東が中央人民政府主席（現在の国家主席）、実務と外交の天才・\aw{周恩来}が政務院総理（現在の国務院総理）となった。

建国当初（1949年9月の臨時憲法である共同綱領）は、いきなり社会主義を目指すのではなく、労働者・農民・小ブルジョワ・民族資本家の四階級の連合による“\aw{新民主主義}（人民民主専政）”の段階とされた。1950年には土地改革法が公布され、何千年も続いた地主制度が暴力的に廃止され、富の公正な分配のため農民に土地が分配された。

\subsection{台湾（中華民国）の独裁と“省籍矛盾”}
一方、国共内戦に敗れて台湾へ逃れた国民党の\aw{\UTF{8523}介石}は、1947年に発効した中華民国憲法の下で、総統として絶対的な独裁体制を敷いた。
台湾には、戦前（日本統治時代を含む）から住んでいた人々（本省人・先住民族）がおり、彼らは大陸から逃げてきた腐敗した国民党政府（外省人）の支配を決して望んでいなかった。この「本省人」と「外省人」の間の深い対立（“\aw{省籍矛盾}”）が爆発したのが、1947年2月に起きた\aw{二・二八事件}である。闇タバコを売る老婦人への警官の暴行を機に全島で民衆が蜂起したが、大陸から派遣された国民党軍により知識人ら2万〜3万人が無差別に虐殺された。

\begin{hocho}[世界最長の「戒厳令」と白色テロ]
二・二八事件の後、\UTF{8523}介石は台湾全土に戒厳令を布告した。この戒厳令は、1949年から\UTF{8523}経国（\UTF{8523}介石の息子）が解除する1987年まで、実に「38年間」も続く世界最長記録の戒厳令となった。この間、共産党のスパイや民主化要求の疑いをかけられた本省人の知識人や学生が、秘密警察によって深夜に連行され、拷問され、処刑される「白色テロ」が日常的に行われた。台湾の目覚ましい経済発展の裏には、この長きにわたる恐怖政治による沈黙の時代があったのである。
\end{hocho}

冷戦構造の変化（1971年の国連代表権喪失など）を経て、\UTF{8523}経国の死後、本省人出身の李登輝が総統に就任。1996年に初めて直接選挙による総統選出が行われ、台湾は民主化を達成した。2000年には民進党の陳水扁が初の政権交代を実現し、以降、国民党と民進党による民主的な選挙が定着している。

% ======================================================================
\section{大躍進の狂気から文化大革命、そして改革開放へ}
% ======================================================================

\subsection{毛沢東の絶対化と大躍進の挫折}
新民主主義の段階を早期に終え、中国は1953年からの\aw{第一次五カ年計画}でソ連の莫大な技術・経済援助を受け、重工業化と農業の集団化（社会主義化）を進めた。
しかし、1956年にソ連のフルシチョフが「スターリン批判」とアメリカとの平和共存路線を打ち出すと、これを「修正主義」と激しく批判した毛沢東との間で\aw{中ソ対立}が表面化。激怒したソ連は技術者を全面撤収させた。

孤立した毛沢東は、ソ連の援助なしに自力更生で一気に共産主義社会へ到達しようと、1958年から\aw{第二次五カ年計画}、いわゆる「\aw{大躍進運動}」を発動した。全国の農村を巨大な\aw{人民公社}に組織化し、鉄鋼と農業の超現実的なノルマを課した。農民たちは鍋や農具を溶かして素人が庭先の「\aw{土法高炉}」で使い物にならないクズ鉄を作り、労働力不足と非科学的な密植農業により農作物は枯死した。結果、推計3000万〜4500万人が餓死するという人類史上最悪の人災を引き起こし、毛沢東は責任を取って国家主席を辞任せざるを得なくなった。

※大躍進の裏で、1959年には中国の抑圧的政策に反発した\aw{チベット反乱}が勃発。ダライ・ラマ14世がインドへ亡命し、インドがこれを匿ったことが、1962年の\aw{中印国境紛争}の直接的な火種となった。

\subsection{プロレタリア文化大革命（文革）の嵐}
大躍進の尻拭いを任された\aw{劉少奇}（国家主席）と\UTF{9127}小平らは、農民の自留地（私有地）や市場原理を一部認める現実的な政策（調整政策）で経済を立て直した。
しかし、権力奪還を画策する毛沢東は、彼らを「資本主義の道を歩む実権派（走資派）」と批判。1966年、若き学生や生徒たちを「\aw{紅衛兵}」として組織し、「造反有理（謀反には理がある）」の狂気のスローガンの下、\aw{プロレタリア文化大革命}を発動した。
紅衛兵たちは党幹部や知識人、教師を市中引き回しにして拷問・殺害し、古い文化財（孔子の墓など）を徹底的に破壊した。劉少奇は凄惨なリンチの末に獄死し、\UTF{9127}小平も失脚してトラクター工場で労働させられた（下放）。国全体が10年間にわたり機能停止に陥る未曾有の大混乱となった。

1969年頃に軍の介入で紅衛兵の暴走が鎮圧されると、毛沢東の絶対的な後継者とされた\aw{林彪}（りんぴょう）が台頭した。しかし毛沢東と対立し、1971年にクーデタを企図して失敗。ソ連へ逃亡する途中でモンゴル上空にて謎の墜落死を遂げた。以後は毛沢東の妻・江青ら急進左派の「\aw{四人組}」（江青、張春橋、姚文元、王洪文）が権力を握り、現実派の周恩来や復活した\UTF{9127}小平と激しい権力闘争を繰り広げた。

\begin{hocho}[\UTF{9127}小平の「白猫黒猫論」とプラグマティズム]
毛沢東が「社会主義的イデオロギー（赤さ）」を絶対視したのに対し、\UTF{9127}小平は「白い猫でも黒い猫でも、ネズミを捕るのが良い猫だ」という有名な「白猫黒猫論」を説いた。つまり、資本主義の要素（黒猫）であれ社会主義（白猫）であれ、人民の生活を豊かにし経済を発展させる政策こそが正しいという極めて現実的（プラグマティズム）な考え方である。この思想が文革期の毛沢東には許しがたい資本主義の毒と映り、\UTF{9127}小平は何度も失脚させられた。しかし、不屈の彼は毛の死後に見事に蘇り、この猫論で13億人を資本主義の海へと導くことになる。
\end{hocho}

\subsection{毛沢東の死と\UTF{9127}小平の「改革開放」}
1976年1月に民衆に慕われた周恩来が死去すると、4月の清明節に天安門広場に集まった追悼の群衆を四人組が武力で弾圧する\aw{第一次天安門事件}が発生した。
同年9月に毛沢東が死去すると、党主席を継いだ\aw{華国鋒}らが軍部と結託して四人組を電撃逮捕し、文革はようやく終結した。
その後、華国鋒を権力闘争で追放した\aw{\UTF{9127}小平}が実権を掌握し、1978年から歴史的な\aw{改革開放政策}を開始。「\aw{四つの現代化}」（工業・農業・国防・科学技術）を掲げ、人民公社を解体して農家の生産請負制を導入し、沿岸部に経済特区（深\UTF{5733}など）を設けて外国資本を導入した。

しかし、経済の自由化は「政治の民主化」の要求を招いた。1989年、民主化に寛容だった胡耀邦（元党総書記）の急死を悼む学生たちが天安門広場を占拠して民主化を要求。これに対し\UTF{9127}小平は、共産党一党独裁の崩壊を恐れ、6月4日に人民解放軍の戦車と銃撃で学生らを無差別に轢き殺す\aw{第二次天安門事件}を引き起こした。
国際社会から激しい非難と制裁を受けたが、同年の冷戦終結（マルタ会談）と1991年のソ連崩壊という歴史的激動（「濤」のごとき時代）を見た\UTF{9127}小平は、1992年に南巡講話を行い「政治の独裁は維持しつつ、経済の資本主義化はさらに加速させる」という奇跡的なアクロバット「\aw{社会主義市場経済}」体制を正式に導入した。この路線は江沢民、胡錦濤へと引き継がれ、そして現在の\aw{習近平}（党総書記が国家主席を兼任する体制）に至り、世界第2位の経済大国にして強力な監視社会を築き上げているのである。

% ======================================================================
\section{国民国家のジレンマと「開発独裁」の台頭}
% ======================================================================

\subsection{「国民の自決」がはらむ排除の論理}
近代ヨーロッパで生まれた「国民国家（nation-state）」と、ウィルソンが提唱した「国民の自決（領域国家による民族自決）」という理念は、独立を目指すアジア・アフリカの人々にとって強力な武器となった。しかし、いざ独立を達成すると、「国民」という概念自体のあいまいさが致命的なジレンマを引き起こした。
西欧諸国が人為的に引いた国境線の中には、言語も宗教も異なる多様な集団が混在していた。一つの「国民」という鋳型に彼らを押し込めようとすると、「国民の名による専制」が発生し、「我々こそが正当な国民である」と主張する多数派が、異質な少数民族や他宗教の集団（「やつら」）を差別し、排斥する構造を生み出したのである。

\subsection{インドにおけるナショナリズムの相克}
この矛盾が最も先鋭的に現れたのがインドである。
ガンディーは南アフリカでの弁護士時代、言語も階級も異なる多様な人々が白人から単に「インド人」と一括されて差別される現実を目の当たりにし、逆説的に「インド人」としての統一の必要性を悟った。彼は\aw{国民会議派}を率い、ヒンドゥー教徒やムスリムなど諸宗教コミュニティの寛容性と「協同」による運動を主張した。
しかし、全インド・ムスリム連盟を率いる\aw{ジンナー}は、多数派ヒンドゥーによる支配を恐れ、当初の協調路線から「ムスリム国家としての分離」を強硬に主張し、結果としてパキスタンの分離と、100万人以上の犠牲を出した大虐殺を引き起こした（さらにパキスタン内部でも、言語の違いから東パキスタンが\aw{バングラデシュ}として独立することになる）。

\begin{hocho}[ガンディーとアンベードカルの死闘、そして仏教改宗]
国民国家の包摂と排除の矛盾を最も鋭く突いたのが、不可触民（ダリット）出身でインド憲法草案起草者となった\aw{アンベードカル}である。ガンディーはカースト制度そのものは否定せず、最下層民を「ハリジャン（神の子）」と呼んで同情的に救済しようとした。しかしアンベードカルは「同情などいらない、法的権利（不可触民の分離選挙枠など）と完全な平等が必要だ」と主張し、ガンディーと激しく対立した。ヒンドゥー教の枠内にいる限り差別は消えないと絶望したアンベードカルは、死の直前の1956年、約50万人の不可触民とともにヒンドゥー教を捨て「仏教」へと集団改宗した。彼の戦いは、多数派のナショナリズムが覆い隠そうとする「内部の差別」への命がけの抵抗であった。
\end{hocho}

\subsection{民主主義をバイパスした「開発独裁」}
独立を果たした新興諸国の至上命題は、貧困からの脱却、すなわち「経済発展と近代化の推進」であった。しかし、民主的な議会政治手続きを踏んでいては、複数のエスニック集団の利害調整に時間がかかりすぎ、国家統合が瓦解する恐れがあった。
そこで、強固な軍事力や強権的な一党支配によって議会政治をバイパスし、上からの強権で強引に資本蓄積と工業化（植民地経済構造からの転換）を推し進める体制が多くの国で誕生した。これを「\aw{開発独裁}」と呼ぶ。
冷戦下のアメリカは、これらの政権が人権を弾圧する独裁であっても、「反共産主義」であり外資に門戸を開く限り、熱烈に支援した。

\begin{itemize}
    \item \textbf{韓国}：\aw{朴正煕}（パク・チョンヒ）政権は1972年の「\aw{10月維新}」で終身大統領制を敷き、国家保安法とKCIAで反対派を弾圧しながら輸出主導型の工業化を達成した。1979年に彼が暗殺された後も、\aw{全斗煥}（チョン・ドゥファン）がクーデタで政権を握り、1980年の\aw{光州事件}で民主化を求める市民を軍で虐殺し、軍事独裁を継続した。
    \item \textbf{インドネシア}：建国の父スカルノが1965年の\aw{九・三〇事件}（国軍左派によるクーデタ未遂）の責任を問われて失脚。この事件を利用して共産党員を推定50万人以上虐殺して実権を握った\aw{スハルト}は、「新体制」と称する開発独裁を30年以上にわたり敷いた。
    \item \textbf{フィリピン}：\aw{マルコス}政権が戒厳令を敷き、政敵アキノを暗殺するなど腐敗した一族独裁を長年継続した（1986年のエドゥサ革命で崩壊）。
    \item \textbf{イラン}：\aw{パフラヴィー2世}が秘密警察サヴァクを用いて反対派を弾圧しつつ、「白色革命」と呼ばれる近代化を強行した。
\end{itemize}

% ======================================================================
\section{ユーゴスラヴィアの「モザイク」と終わらぬ中東紛争}
% ======================================================================

\subsection{ティトーとユーゴスラヴィアの独自の道}
\begin{center}
  \includegraphics[width=10cm]{yugo.png}
\end{center}
バルカン半島に位置するユーゴスラヴィア（1918年建国、1929年改称）は、第二次大戦中、共産党の\aw{ティトー}率いるパルチザンがソ連軍に頼らず「\aw{自力でドイツ軍を撃退}」したという強い誇りを持っていた。
戦後、ソ連主導で結成された\aw{コミンフォルム}（当初本部はユーゴのベオグラード）に加盟したが、スターリンの画一的な支配にティトーが反発。1948年に「民族主義的」「非協調的」であるとしてコミンフォルムから除名された。
孤立したユーゴはソ連の真似をやめ、工場や農場の運営を国家ではなく労働者自身に任せる「\aw{自主管理社会主義}」という独自路線を打ち出した。外交面では、米ソのどちらにも属さない\aw{非同盟政策}（積極的平和共存）をとり、第三世界運動の旗手となった。
しかし、この国は「6つの共和国、5つの民族、4つの言語、3つの宗教、2つの文字、1つの国家」と呼ばれる極めて複雑なモザイク国家であった。ティトーという絶対的なカリスマの死（1980年）後、経済危機の進行とともに各共和国の民族主義が暴走。1991年にスロヴェニアとクロアティアが独立を宣言したのを皮切りに、凄惨な民族浄化を伴う\aw{ユーゴスラヴィア紛争}へと転落し、国家は完全に解体された。

\begin{hocho}[モザイク国家]
  ティトー大統領のもとで結束したユーゴスヴィアは、極めて複雑に民族や文化の入り組んだ国家であった。\\
6つの共和国：スロヴェニア、クロアティア、ボスニア・ヘルツェゴビナ、
セルビア、モンテネグロ、マケドニア\\
5つの民族：スロヴェニア、クロアティア、セルビア、モンテネグロ、マケドニア\\
4つの言語：セルビア語、スロヴェニア語、クロアティア語、マケドニア語\\
3つの宗教：カトリック、東方正教、イスラーム\\
2つの文字：ラテン文字、キリル文字\\
1つの国家：ユーゴスラヴィア社会主義連邦共和国
\end{hocho}


\subsection{パレスチナ問題と四次の中東戦争}
中東最大の火薬庫であるパレスチナ問題は、イギリスが二枚舌外交の処理を諦めて国際連合に丸投げしたことから始まる。
1947年、国連はパレスチナ分割案と\aw{イェルサレムの国際管理地域化}を決議。1948年5月、ユダヤ人が\aw{イスラエル}建国を宣言すると、これを認めない周辺アラブ諸国が侵攻し、\aw{パレスチナ戦争}（第1次中東戦争、1948-49年）が勃発した。アメリカとイギリスの援助を受けたイスラエルが勝利して領土を拡大し、故郷を追われた100万人以上の「パレスチナ難民」が発生した。

\begin{itemize}
    \item \aw{第3次中東戦争}（1967年）：イスラエルがエジプト・シリア・ヨルダンに先制の電撃攻撃を行い、わずか6日間でシナイ半島、ヨルダン川西岸、ゴラン高原などを占領。アラブ側に圧倒的な挫折感を与えた。
    \item \aw{第4次中東戦争}（1973年）：失地回復を狙うエジプトとシリアが、ユダヤ教の休日に奇襲をかけて勃発。この時、アラブ石油輸出国機構（OAPEC）が禁輸と減産という「石油戦略」を発動し、世界経済を直撃する第一次石油危機（オイルショック）を引き起こした。
\end{itemize}

\begin{multicols}{2}
\subsection{和平の模索と挫折}
四度にわたる戦争の疲弊から、エジプトの\aw{サダト}大統領はイスラエルのベギン首相と和平交渉を行い、1978年に米カーター大統領の仲介で\aw{キャンプ＝デーヴィッド合意}に達し、1979年に平和条約を調印した。
一方、1964年に結成された\aw{パレスチナ解放機構（PLO）}は、当初はイスラエル打倒を掲げてテロ戦術（ミュンヘン五輪事件など）を展開していた。しかし、1987年にイスラエル占領地でパレスチナ民衆の投石による抵抗運動（インティファーダ）が激化すると、1988年にPLOの\aw{アラファト}議長はイスラエルとの共存とパレスチナ国家樹立へ方針を転換。1993年、米クリントン大統領の仲介でイスラエルの\aw{ラビン}首相との間で\aw{オスロ合意}（パレスチナ暫定自治協定）に調印し、自治政府が活動を始めた。しかし、ラビン首相の暗殺やユダヤ人入植地の拡大などにより和平は頓挫し、現在に至るまで流血の惨事が続いている。\\\\\\\\\\\\\\\\\\


\begin{center}
  
  \includegraphics[width=7cm]{israelimap.png.png}
\end{center}
\end{multicols}

% ======================================================================
\section{中東の革命：イラクとイランの激動}
% ======================================================================

\subsection{イラク革命による王政崩壊}
第一次世界大戦後、イギリスはメッカの守護者であったハーシム家（フサインの息子ファイサル）をイラク国王に据え、親英的な統治を行っていた。しかし、冷戦下でイギリスが主導する反共軍事同盟（バグダード条約機構）に組み込まれたことに対し、アラブ民族主義に燃える軍部が反発。1958年に将校団によるクーデタ（\aw{イラク革命}）が起き、王一族は惨殺され共和政へと移行した。当然、バグダード条約機構からも脱退した。

\subsection{イラン・イスラーム革命の衝撃}
イランでは、親米・親英の独裁体制を敷く\aw{パフラヴィー2世}に対し、1951年に民族主義者のモサデグ首相がイギリス資本に独占されていた「石油の国有化」を断行した。これに激怒した英米は、CIAの秘密工作（クーデタ）によってモサデグ政権を転覆させ、国王の権力を復活させた。
帰国した国王は、アメリカの莫大な援助と石油収入を背景に、農地改革や女性参政権の付与など、上からの強引な近代化政策「\aw{白色革命}」を1963年から推進した。しかし、これによって貧富の格差は異常に拡大し、欧米の退廃的な文化の流入は敬虔なイスラーム教徒の激しい怒りを買った。

1979年、国王の秘密警察による拷問と弾圧に耐えかねた民衆が蜂起。パリに亡命していた厳格なイスラーム法学者\aw{ホメイニ}の指導のもと革命が成功し、国王は国外へ逃亡。親米のパフラヴィー朝は崩壊し、\aw{イラン＝イスラーム共和国}が成立した（\aw{イラン革命}）。
この新国家は、シーア派の論理に基づく「法学者の統治論」を採用し、コーランと預言者の言行録（\aw{ハディース}）に基づくイスラーム法（\aw{シャリーア}）を絶対視した。「欧米的なもの」を徹底的に排除し、アメリカ大使館人質事件を起こすなど、反米を強烈な国是とした。

\begin{hocho}[イマームの代行者とイラン・イラク戦争]
スンナ派が代々の「カリフ」を指導者とするのに対し、イランの国教であるシーア派は、預言者ムハンマドの血を引く「\aw{イマーム}」のみを真の指導者とみなす。現在イマームは姿を隠している（お隠れ）ため、最高指導者（ホメイニら）はその代理人として絶対的な権威を持った。\\
このイスラーム革命の熱狂が、自国のシーア派住民に波及することを極度に恐れたのが、隣国イラクの世俗主義的な独裁者サダム・フセインである。イラクは1980年にイランに侵攻し、泥沼の\aw{イラン＝イラク戦争}（1980-88年）が勃発した。アメリカは「反米のイランを叩くため」にイラクのフセイン政権に化学兵器の原料を含む裏の支援を行い、これが後に湾岸戦争やイラク戦争の悲劇へと繋がるのである。
\end{hocho}

% ======================================================================
\section{「第三世界」運動の連帯と挫折、そして欧州統合}
% ======================================================================

\subsection{非同盟主義の台頭}
冷戦による米ソの核戦争の恐怖が迫る中、「西（資本主義陣営）でも東（社会主義陣営）でもない」新たな勢力として、アジア・アフリカの新興諸国が結集し「\aw{第三世界}」を形成した。この\aw{非同盟主義}は、民族運動の強力な先導役となった。
\begin{itemize}
    \item \textbf{平和五原則}（1954年）：中国の\aw{周恩来}総理とインドの\aw{ネルー}首相が共同声明として発表。（領土と主権の相互尊重、相互不可侵、内政不干渉、平等互恵、平和共存）。
    \item \textbf{第1回アジア・アフリカ会議（バンドン会議）}（1955年）：インドネシアのバンドンで開催。インド、セイロン、インドネシア（スカルノ）、ビルマ、パキスタンの5カ国の首相が提唱。29カ国が参加し、反植民地主義と平和五原則を拡張した「\aw{平和十原則}」を宣言した。
    \item \textbf{非同盟諸国首脳会議}（1961年）：アジア・アフリカ諸国に加え、欧州からユーゴスラヴィアが参加し、第1回会議がベオグラードで開催された。ティトー（ユーゴ）、ネルー、ナセル（エジプト）、スカルノ、ンクルマ（ガーナ）らが第三世界の旗手として「平和共存」と「反植民地主義」を強烈にアピールした。
\end{itemize}

\begin{hocho}[周恩来の神業外交とバンドン会議]
バンドン会議には、親米（フィリピン等）、親ソ（中国等）、非同盟（インド等）の様々な立場の国が参加しており、当初はイデオロギー対立で紛糾した。とくに「ソ連の東欧支配も新たな植民地主義ではないか」という批判が出た時、中国の周恩来は絶妙な機転を利かせた。彼は用意していた強硬な演説原稿を捨て、「我々は共通の基盤を求めてここに来たのであり、争いを求めて来たのではない」と穏やかにスピーチし、アメリカへの名指しの批判も避けた。この周恩来の寛容でスマートな態度は各国の警戒心を解き、会議を大成功（平和十原則の採択）に導いたのである。
\end{hocho}

\subsection{第三世界運動の退潮}
しかし、1960年代以降、華々しかった第三世界の結束は急速に崩壊していった。担い手であった指導者たちが影響力を失ったためである。
1959年のチベット動乱をきっかけに、平和五原則を共に提唱した中国とインドが軍事衝突（1962年の\aw{中印国境紛争}）。1965年にはバンドンのホストであったインドネシアのスカルノが\aw{九・三〇事件}で失脚。1967年には非同盟の顔であったエジプト（ナセル）らアラブ諸国が\aw{第3次中東戦争}でイスラエルに惨敗。各国の内部紛争や経済的困窮が表面化し、連帯の夢は後退した。

\subsection{ヨーロッパ統合の歩み}
第三世界が退潮する一方で、二度の世界大戦の惨禍を経験した西ヨーロッパでは、「軍需産業の基盤である資源を共同管理し、戦争を物理的に不可能にする」という統合の試みが着実に進展していた。
\begin{itemize}
    \item 1950年のシューマン・プラン（仏外相の提唱）を機に、1952年に\aw{ヨーロッパ石炭鉄鋼共同体（ECSC）}が発足。
    \item 1957年の\aw{ローマ条約}により、仏・西独・伊・ベネルクス三国によって1958年に\aw{ヨーロッパ経済共同体（EEC）}とヨーロッパ原子力共同体（EURATOM）が発足。
    \item 1967年、これら3つの組織が統合され\aw{ヨーロッパ共同体（EC）}となった（イギリスは当初EFTAを作り対抗したが、1973年にECへ加盟）。
    \item 冷戦終結の熱狂の中、1992年に\aw{マーストリヒト条約}が調印され、1993年に政治的・通貨的統合を進める\aw{ヨーロッパ連合（EU）}が発足した。
\end{itemize}
その後、EUは中・東欧の旧社会主義国を多数加盟させ東方へ拡大したが、移民問題や経済格差から新たな対立が生じ、2016年の国民投票を経て2020年にイギリスが離脱する（\aw{Brexit}）という歴史的逆流も発生している。

\begin{center}
  \includegraphics[width=15cm]{eruop.png}
\end{center}

% ======================================================================
\section{ラテンアメリカと新興諸国の苦境}
% ======================================================================

\subsection{アメリカの「裏庭」と軍事政権の時代}
ラテンアメリカは長らく「アメリカの裏庭」として扱われてきた。戦後の冷戦下、アメリカは1947年にパン＝アメリカ会議でリオ協定を結び、1948年に\aw{米州機構（OAS）}を発足させて共産主義の封じ込めを図った。
しかし、1959年1月にカストロやチェ・ゲバラらが親米独裁政権を打倒する\aw{キューバ革命}を成功させ、その後社会主義を宣言すると、アメリカは喉元に匕首を突きつけられた恐怖を抱いた（1966年にハバナで反帝国主義の会議が開かれるなど、革命の輸出拠点となった）。

このキューバ革命の波及を防ぐため、アメリカ（CIA）は1960年代後半から、ラテンアメリカ各国の軍部を裏で支援し、左派的・民族主義的な政権を次々と暴力的に転覆させた。
\begin{itemize}
    \item \textbf{チリ}：1970年に世界で初めて民主的選挙で成立したアジェンデ左派政権に対し、1973年に軍部のピノチェトが軍事クーデタを起こし、大統領府を爆撃してアジェンデを殺害。以後、過酷な軍事独裁を敷いた。
    \item \textbf{アルゼンチン}：1976年に親米軍事政権が成立。数万人の反体制派市民が秘密裏に誘拐・殺害（死のフライト）された。不満を逸らすために1982年にイギリスと\aw{フォークランド戦争}を起こすも惨敗し、軍政は崩壊して民政移管した。
\end{itemize}
アメリカは、自国の多国籍企業や大土地所有者の利権を守ってくれる政権であれば、いかに残虐な独裁政権であっても支持し、必要であれば直接軍事介入（1989年のパナマ侵攻など）も辞さなかった。1970年代のラテンアメリカは、恐怖と血に塗れた\aw{軍事政権（開発独裁）}の時代であった。

\begin{hocho}[シカゴ・ボーイズとチリの「人体実験」]
チリのピノチェト軍事政権は、反対派をスタジアムに集めて虐殺する一方で、経済政策についてはアメリカのシカゴ大学でミルトン・フリードマンに学んだ新自由主義経済学者のグループ、通称「シカゴ・ボーイズ」に丸投げした。彼らは、医療や教育、年金に至るまで国家の福祉を徹底的に削減し、あらゆるものを民営化する過激な自由市場政策をチリに導入した。この「軍事力で国民の抵抗を完全に封じ込めた上で行われた究極の資本主義の人体実験」は、マクロ経済の数値を改善させた一方で、富裕層への莫大な富の集中と、絶望的な貧困層を生み出すこととなった。
\end{hocho}

\begin{center}
  \includegraphics[width=15cm]{same.png}
\end{center}

\subsection{「南北問題」と従属理論}
独立した新興国の最大の苦境は、植民地時代から続く歪な経済構造にあった。外貨を獲得するためには、特定の一次産品（コーヒー、カカオ、銅など）の輸出に依存せざるを得ない「\aw{モノカルチャー化}」から抜け出せず、国際価格の暴落によって国家経済が容易に破綻した。
1960年代、豊かな北の先進工業国と、貧困にあえぐ南の途上国との圧倒的な経済格差が「\aw{南北問題}」として国際社会で提起された。

欧米の学者は「いずれ近代化する」と楽観視したが、ラテンアメリカの経済学者らはこれを痛烈に批判した。彼らが唱えた「\aw{従属理論}」によれば、南の貧困は単なる発展の遅れではなく、「500年来の歴史が産み出した、世界経済の中心（北）が周辺（南）を恒常的に搾取し続ける構造そのものに原因がある」とされた。

\subsection{NICsの台頭と南南問題}
1970年代以降、人件費の高騰した先進国から、安い労働力と規制の緩さを求めて製造工場が途上国へ移転し始めた。強権的な開発独裁の下で労働運動を弾圧しつつ、この外資を積極的に導入して「輸出主導型」の工業化を進め、急激に経済成長を遂げた国々が現れた。
これらを\aw{新興工業諸国（NICs）}または新興工業経済地域（NIEs）と呼ぶ。とくにアジアでは\aw{香港}、\aw{シンガポール}、\aw{大韓民国}、\aw{台湾}が「\aw{アジア四小龍}」と呼ばれ急成長し、中南米でもメキシコやブラジルが台頭した。
しかし、この一部の国の台頭は、未だ一次産品輸出から抜け出せない最貧国との間に、第三世界内部での新たな格差拡大（\aw{南南問題}）を生み出す結果となった。

% ======================================================================
\section{冷戦後の世界：終わらない紛争と新興大国の台頭}
% ======================================================================

\subsection{エスニック紛争と難民の世紀}
冷戦が終結し、米ソの超大国による「抑え付け」が消滅すると、新興国内部でくすぶっていた政治的・経済的な利益配分の不平等が爆発した。複数の集団を抱え込む国家において、少数派としてのアイデンティティ（\aw{エスニシティ}）を核とした分離独立運動や、政権をめぐる凄惨な内戦が頻発した。
冷戦後の紛争は国家間の戦争ではなく、国家内の「エスニック紛争（民族浄化）」が主流となり、非戦闘員への凄惨な人権侵害により大量の難民が発生した。これに対処するため、国連平和維持活動（\aw{PKO}）や\aw{国連難民高等弁務官事務所（UNHCR）}、そして「国境なき医師団」などの民間組織（\aw{NGOs}）の果たす役割がかつてなく増大している。

\subsection{深刻化する地域紛争}
\begin{itemize}
    \item \textbf{カシミール問題}：インド・パキスタン戦争の原因となったカシミール地方は、現在もインド支配地域、パキスタン支配地域、中国支配地域に分断されており、両国が核保有国となった現在、世界で最も危険な火薬庫の一つとなっている。
    \item \textbf{ソマリアの崩壊}：アフリカの角では、冷戦下でソマリアとエチオピアが\aw{オガデン戦争}（1977-88年）を戦い、敗北したソマリアは国内の軍閥が対立して\aw{ソマリア内戦}（1991年〜）へと転落。国家が事実上崩壊し、海賊の跳梁やイスラーム過激派のテロの温床となっている（2012年に統一政府が成立するも依然不安定）。
\end{itemize}

\begin{hocho}[ブラックホーク・ダウンとアメリカのトラウマ]
1993年、無政府状態のソマリアで人道支援物資の略奪を繰り返す軍閥を捕らえるため、アメリカ軍の特殊部隊が首都モガディシュに降下した。しかし、作戦中にハイテク・ヘリコプター「ブラックホーク」が民兵のRPG（ロケット弾）で撃墜され、救出に向かった米兵たちが数万人の武装民衆に完全に包囲されるという悪夢の市街戦（モガディシュの戦闘）が発生した。米兵の遺体が群衆に引きずり回される映像が全世界に流れ、クリントン大統領は部隊の完全撤退を余儀なくされた。このトラウマにより、アメリカは翌1994年のルワンダ大虐殺（約80万人が犠牲）への軍事介入を躊躇することになり、冷戦後の「世界の警察官」としての役割に限界を示した。
\end{hocho}

\subsection{新興経済諸国（BRICS）の台頭と新冷戦}
冷戦後のグローバル経済の中で、巨大な人口や豊富な資源を背景に急成長を遂げた大国群を\aw{BRICS}（ブラジル、ロシア、インド、中国、南アフリカ）と呼ぶ。
\begin{itemize}
    \item \textbf{ロシア}：ソ連崩壊後、\aw{エリツィン}大統領が急速な市場経済導入を行って経済が崩壊状態に陥ったが、1999年に登場した\aw{プーチン}は、豊富なエネルギー資源を背景に強権で国家を立て直した。近年は欧米への対抗姿勢を鮮明にし、ウクライナ侵攻など「\aw{新冷戦}」の震源地となっている。
    \item \textbf{中国}：\UTF{9127}小平の「社会主義市場経済」路線を継承した江沢民、胡錦濤時代に世界第2位の経済大国へ躍進。続く\aw{習近平}（2012年〜）は、“中華民族の偉大な復興”を掲げて圧倒的な大国姿勢（戦狼外交）を示し、アメリカの覇権に真っ向から挑戦している。
    \item \textbf{ベトナム}：中国に倣い、社会主義の政治体制を維持したまま1986年に市場経済を導入する\aw{ドイモイ（刷新）}政策を採用し、外国企業を誘致して著しい経済発展を遂げている。
\end{itemize}
このように世界は多極化の時代を迎えたが、経済発展をとげた新興国であっても、国内に深刻な経済格差や抑圧された民族問題を抱えており、21世紀の国際社会は極めて不安定な状況を内包したまま進んでいるのである。
% ======================================================================
\clearpage
\thispagestyle{fancy}
\fancyhead[L]{\textbf{世界史探究 一問一答（全100問）}}
\fancyhead[R]{\textbf{戦後のアジア・第三世界}}

\section*{【総復習】世界史現代史 一問一答チェック（全100問）}
\addcontentsline{toc}{section}{【総復習】一問一答チェック}

\footnotesize % さらにページに詰め込むため文字サイズを少し縮小
\begin{multicols}{2} % 2段組み開始

% 一問一答用専用マクロ（常に赤字出力）
\newcounter{qcounter}
\setcounter{qcounter}{0}
\newcommand{\QA}[2]{%
  \stepcounter{qcounter}%
  \noindent\textbf{Q\theqcounter.} #1 \dotfill \textcolor{red}{\textbf{#2}}\par\vspace{1.2mm}%
}

\subsection*{【アジア・中東・アフリカの独立と混乱】}
\QA{第二次大戦後、イギリス連邦が名称を変更した、より対等な主権国家の連合を何というか。}{コモンウェルス・オブ・ネイションズ}
\QA{1948年に共和国法を可決し、1949年にイギリス連邦から完全に脱退した国はどこか。}{アイルランド}
\QA{フランスが第5共和制下で形成し、現在はフランス語圏の文化的紐帯となっている組織は何か。}{フランコフォニー}
\QA{1992年に社会主義を放棄し「モンゴル人民共和国」から国名を改称した国はどこか。}{モンゴル国}
\QA{オランダの再支配に抵抗して独立を勝ち取ったインドネシアの初代大統領は誰か。}{スカルノ}
\QA{1946年、戦前からの約束通りアメリカから独立を果たした東南アジアの国はどこか。}{フィリピン}
\QA{1965年にマレーシア連邦から分離独立を余儀なくされた、華人主体の国はどこか。}{シンガポール}
\QA{シンガポールの初代首相で、強権的な手法で国を急成長させた人物は誰か。}{リー＝クアンユー}
\QA{ビルマ（現ミャンマー）の独立運動を指導した「建国の父」は誰か。}{アウン・サン}
\QA{インドとパキスタンが帰属を巡って三度の戦争を引き起こした地方はどこか。}{カシミール地方}
\QA{1971年、西パキスタンの支配に反発し、インドの支援を受けて分離独立した国はどこか。}{バングラデシュ}
\QA{バングラデシュ独立の中心となり、現在も同国の首都である都市はどこか。}{ダッカ}
\QA{カースト差別撤廃を巡りガンディーと激しく対立した、インド憲法草案の起草者は誰か。}{アンベードカル}
\QA{フランス委任統治領から1943年に独立し、のちに凄惨な宗教内戦を経験した国はどこか。}{レバノン}
\QA{1958年にエジプトとシリアが合邦して成立した国家（61年に解消）を何というか。}{アラブ連合共和国}
\QA{1952年にクーデタを起こし、エジプトの王政を打倒した軍内の秘密組織を何というか。}{自由将校団}
\QA{アスワン＝ハイダム建設資金確保のため、1956年にスエズ運河国有化を宣言した大統領は誰か。}{ナセル}
\QA{スエズ運河国有化に反発し、英仏とイスラエルがエジプトを攻撃した戦争を何というか。}{スエズ戦争}
\QA{イラクで1958年にクーデタを起こして王政を打倒した革命を何というか。}{イラク革命}
\QA{イラク革命後、イラクが脱退した親西側の反共軍事同盟は何か。}{バグダード条約機構}
\QA{フランス本国を巻き込む泥沼の独立戦争を展開した、アルジェリアの組織（略称）は何か。}{FLN}
\QA{1962年、ド＝ゴール政権下のフランスがアルジェリアの独立を承認した協定を何というか。}{エビアン協定}
\QA{1957年、サハラ以南で最初に独立したガーナの初代指導者は誰か。}{ンクルマ}
\QA{フランス領を中心にアフリカ諸国17カ国が一挙に独立を果たした1960年は何と呼ばれたか。}{アフリカの年}
\QA{1963年、パン＝アフリカニズムの理念の下に設立されたアフリカ諸国の連帯組織の略称は何か。}{OAU}
\QA{モロッコが実効支配し、ポリサリオ戦線が独立を求めて武装闘争を行っている地域はどこか。}{西サハラ}
\QA{アンゴラ領の一部であるが、コンゴ民主共和国を挟んで飛び地となっている地域はどこか。}{カビンダ}
\QA{1968年に独立し、2018年に「エスワティニ王国」へ国名変更したアフリカ南部の国はどこか。}{スワジランド}
\QA{1991年にソマリアから独立を宣言したが、国際的に未承認の国はどこか。}{ソマリランド共和国}
\QA{1974年、本国の軍事独裁が崩壊し、アフリカ植民地の独立につながった革命は何か。}{ポルトガル革命}
\QA{南アフリカ共和国で1991年に撤廃が宣言された、極端な人種隔離政策を何というか。}{アパルトヘイト}
\QA{1990年に南アフリカから独立し、アフリカの脱植民地化がほぼ完了したとされる国はどこか。}{ナミビア}

\subsection*{【中国の激動と台湾の歩み】}
\QA{1949年、国共内戦に勝利して中華人民共和国を建国した中国共産党の指導者は誰か。}{毛沢東}
\QA{中華人民共和国の建国当初に掲げられた、四階級の連合による国家体制を何と呼ぶか。}{新民主主義}
\QA{建国初期の中国で、周恩来とともに国務院総理として実務・外交を担った人物は誰か。}{周恩来}
\QA{台湾に逃れた国民党の\UTF{8523}介石政権に対し、1947年に本省人の不満が爆発し弾圧された事件は何か。}{二・二八事件}
\QA{台湾で世界最長となる38年間敷かれ、言論弾圧や白色テロの根拠となった法令は何か。}{戒厳令}
\QA{1958年から毛沢東が発動し、非科学的な農工業政策により数千万人の餓死者を出した政策は何か。}{大躍進運動}
\QA{大躍進運動の際、農・工・商・学・兵を一体化した農村社会の基盤として組織されたものは何か。}{人民公社}
\QA{中国の社会主義化に反発し、1959年にチベットで起きた反乱を何というか。}{チベット反乱}
\QA{1962年、チベット反乱への対応や国境線を巡って中国とインドの間で起きた武力衝突は何か。}{中印国境紛争}
\QA{毛沢東が大躍進の失敗から権力奪還を狙い、1966年に発動した10年に及ぶ一大政治闘争は何か。}{文化大革命}
\QA{文化大革命で「造反有理」を掲げて破壊活動を行った若者や学生らの組織を何というか。}{紅衛兵}
\QA{文化大革命で毛沢東に「走資派」として打倒・迫害された当時の国家主席は誰か。}{劉少奇}
\QA{文化大革命の最中に毛沢東の後継者とされたが、対立して逃亡中に墜落死した人物は誰か。}{林彪}
\QA{毛沢東の死後、権力を握ろうと画策した江青ら急進左派の4人組を何というか。}{四人組}
\QA{四人組を逮捕して文化大革命を終わらせ、一時的に党主席となった人物は誰か。}{華国鋒}
\QA{文革で失脚したが復活し、1978年から「改革開放政策」を主導した人物は誰か。}{\UTF{9127}小平}
\QA{\UTF{9127}小平が説いた、「資本主義でも社会主義でも、経済を発展させる政策が正しい」とする例え話を何というか。}{白猫黒猫論}
\QA{1989年、民主化を求める学生らを人民解放軍が武力弾圧した事件を何というか。}{第二次天安門事件}
\QA{\UTF{9127}小平が1992年に導入した、共産党の独裁を維持しつつ経済の資本主義化を進める体制を何というか。}{社会主義市場経済}
\QA{現在の中国の最高指導者であり、「中華民族の偉大な復興」というスローガンを掲げる人物は誰か。}{習近平}

\subsection*{【朝鮮戦争とベトナム戦争】}
\QA{1948年に成立した大韓民国の初代大統領で、四月革命により退陣した人物は誰か。}{李承晩}
\QA{日本統治時代の開発により、朝鮮半島の北部で重工業が、南部で農業が盛んであった構造を何と呼ぶか。}{南農北工}
\QA{1950年代の韓国で、アメリカの余剰農作物を利用して発達した製粉・精糖・紡織工業の総称は何か。}{三白産業}
\QA{1950年に北朝鮮が南進して勃発し、米ソの代理戦争となった戦争は何か。}{朝鮮戦争}
\QA{朝鮮戦争の勃発から休戦までの間、事実上の軍事境界線となった緯度は何度か。}{北緯38度線}
\QA{朝鮮戦争において、マッカーサーが総司令官を務め、アメリカ軍を主体として編成された軍隊を何というか。}{国連軍}
\QA{朝鮮戦争において、北朝鮮の危機に際して中国が派遣した大部隊を何というか。}{人民義勇軍}
\QA{1961年に軍事クーデタを起こし、強権のもと韓国の高度経済成長を推進した大統領は誰か。}{朴正煕}
\QA{朴正煕政権が経済協力金を得るため、国内の反対を押し切って1965年に日本と結んだ条約は何か。}{日韓基本条約}
\QA{1980年、全斗煥の軍事クーデタに反対し民主化を求めた市民を軍が武力弾圧した事件を何というか。}{光州事件}
\QA{強権で民主主義を抑圧しつつ、上からの資本蓄積で経済成長を推進する政治体制を何と呼ぶか。}{開発独裁}
\QA{1945年にベトナム民主共和国の独立を宣言した国民的英雄は誰か。}{ホー＝チ＝ミン}
\QA{1954年にベトナム人民軍がフランス軍を大破・降伏させた激戦を何というか。}{ディエンビエンフーの戦い}
\QA{1954年に結ばれ、北緯17度線を境界にベトナムを一時分割した協定は何か。}{ジュネーヴ協定}
\QA{アメリカのアイゼンハウアーらが唱えた、「一国が共産化すれば周辺国も次々と赤化する」という理論は何か。}{ドミノ理論}
\QA{1975年、内戦を経てラオスに成立した社会主義国を何というか。}{ラオス人民民主共和国}
\QA{アメリカの支援を受けて南ベトナムに独裁政権を樹立した大統領は誰か。}{ゴ＝ディン＝ジエム}
\QA{ゴ政権の独裁に反発し、1960年に結成された反政府ゲリラ組織を何というか。}{南ベトナム解放民族戦線}
\QA{1964年、アメリカが北爆を開始する口実とした捏造事件（トンキン湾事件）時の大統領は誰か。}{ジョンソン}
\QA{ベトナム戦争中、アメリカ軍がジャングルを破壊するために大量に散布した化学薬品は何か。}{枯葉剤}
\QA{1968年の旧正月、解放民族戦線が南部の都市で一斉に蜂起し、アメリカの反戦運動を激化させた作戦は何か。}{テト攻勢}
\QA{「戦争のベトナム化」を掲げて米軍撤退を開始する一方、カンボジア等へも侵攻したアメリカ大統領は誰か。}{ニクソン}
\QA{1973年、米軍の完全撤退を取り決めた協定を何というか。}{ベトナム和平協定}

\subsection*{【第三世界と冷戦後の世界・ヨーロッパ統合】}
\QA{1954年、周恩来とネルーが共同声明として発表した平和共存の原則を何というか。}{平和五原則}
\QA{1955年にインドネシアで開催され、「平和十原則」を採択した歴史的な会議は何か。}{バンドン会議}
\QA{ソ連と距離を置き「自主管理社会主義」路線をとったユーゴスラヴィアの指導者は誰か。}{ティトー}
\QA{1961年、東西どちらにも属さない新興国が第1回非同盟諸国首脳会議を開催した都市はどこか。}{ベオグラード}
\QA{1965年、インドネシアでスカルノが失脚しスハルトが実権を掌握する契機となった事件は何か。}{九・三〇事件}
\QA{1948年にユダヤ人が建国を宣言し、周辺アラブ諸国との間で中東戦争を引き起こした国はどこか。}{イスラエル}
\QA{1973年の第4次中東戦争の際、アラブ石油輸出国機構（OAPEC）が発動して世界経済を直撃した政策は何か。}{石油戦略}
\QA{1978年、エジプトのサダト大統領とイスラエルのベギン首相が単独講和に合意した枠組みを何というか。}{キャンプ＝デーヴィッド合意}
\QA{1964年に結成され、アラファト議長のもとイスラエルとの共存路線へ転換した組織の略称は何か。}{PLO}
\QA{1993年、イスラエルのラビン首相とPLOの間で結ばれたパレスチナ暫定自治協定を何というか。}{オスロ合意}
\QA{イランのパフラヴィー2世が進めた、上からの急激な近代化政策を何と呼ぶか。}{白色革命}
\QA{1979年にイラン革命を指導し、反米の「イラン＝イスラーム共和国」を樹立した法学者は誰か。}{ホメイニ}
\QA{冷戦下のラテンアメリカで、アメリカが主導して1948年に結成した反共封じ込めの国際機関（略称）は何か。}{OAS}
\QA{1959年、カストロやチェ・ゲバラらが親米政権を打倒して社会主義国を打ち立てた革命は何か。}{キューバ革命}
\QA{1979年、ニカラグアでソモサ独裁政権を打倒した革命を指導した組織を何というか。}{サンディニスタ民族解放戦線}
\QA{1970年に世界初の選挙による社会主義政権を樹立したが、軍のクーデタで倒れたチリの大統領は誰か。}{アジェンデ}
\QA{1982年、国内の不満をそらすためアルゼンチンの軍事政権がイギリスと戦って大敗した戦争は何か。}{フォークランド戦争}
\QA{豊かな北半球の先進工業国と、貧しい南半球の途上国との圧倒的な経済格差問題を何というか。}{南北問題}
\QA{途上国の貧困は先進国による恒常的な搾取の構造に起因するという、ラテンアメリカ発祥の経済理論は何か。}{従属理論}
\QA{1970年代以降、輸出主導で急速な工業化を達成した韓国や台湾などの「新興工業諸国」の略称は何か。}{NICs}
\QA{NICsの台頭により、途上国間で新たに生じた経済格差問題を何と呼ぶか。}{南南問題}
\QA{1954年、アメリカのビキニ環礁での水爆実験で被曝し、原水爆禁止運動高揚の契機となった日本の漁船は何か。}{第五福竜丸}
\QA{1950年のストックホルム・アピール等を背景に、1957年に世界の科学者が集って開かれた反核会議は何か。}{パグウォッシュ会議}
\QA{1952年、フランス外相の提案により結成した、欧州統合の第一歩となる組織（略称）は何か。}{ECSC}
\QA{ローマ条約に基づき、1958年に発足した「ヨーロッパ経済共同体」の略称は何か。}{EEC}
\QA{1992年に調印され、翌年ヨーロッパ連合（EU）を発足させる根拠となった条約は何か。}{マーストリヒト条約}
\QA{冷戦後、急速に台頭したブラジル、ロシア、インド、中国、南アフリカの新興大国群の総称は何か。}{BRICS}
\QA{ソ連崩壊後のロシアで、急速な市場経済導入（ショック・セラピー）を行い社会の混乱を招いた大統領は誰か。}{エリツィン}
\QA{1986年、ベトナムが市場経済を導入して外資誘致と経済成長を目指した政策を何と呼ぶか。}{ドイモイ}
\QA{国民国家の枠内で、独自の言語や文化を共有する少数派が主張する帰属意識（アイデンティティ）を何というか。}{エスニシティ}
\QA{冷戦後、多発する難民問題に対処するため、その保護や支援活動を牽引している国連機関（略称）は何か。}{UNHCR}
\QA{冷戦後の紛争解決のため派遣される「国連平和維持活動」のアルファベット略称は何か。}{PKO}
\QA{冷戦後、ユーゴスラヴィア等で多発した、国内の民族的帰属意識の違いによる紛争を何というか。}{エスニック紛争}
\QA{オガデン戦争の敗北後に無政府状態に陥り、1991年から激しい内戦が続いているアフリカの国はどこか。}{ソマリア}

\end{multicols}
\normalsize


% ページ合わせコマンド
\end{document}